% s02_setup. Section file; new macros at the top; label prefix: set
\section{Set-up: local \texorpdfstring{$\bbP^2$}{P2}, mirror data and the graph sum}
\label{sec:set}

This section fixes the objects of the paper. The target, the mirror data and the Givental action are recalled from \cite{BS} and \cite{PPZ15}. The rows, the factors and the decoration sum are defined here by explicit formulas, and the statements of \cref{sec:lea,sec:van,sec:val} and of the first half of \cref{sec:reg}, and those of \cref{sec:row} listed at its opening, are statements about these formulas.

%-------------------------------------------------------------------------------
\subsection{Target, tori and invariants}
\label{sec:set-target}

Let $Z\coloneqq \KPtwo\times\bbC^2$, with projection $\mathrm{pr}:Z\to\bbP^2$, and let $\epsilon_1,\epsilon_2$ be the equivariant parameters of the $\bbC^2$ factor. Throughout,
\[
E\coloneqq\epsilon_1+\epsilon_2,\qquad \epsilon_1\epsilon_2E\neq 0.
\]
The letter $E$ with no argument is always this weight; the function of \cite{BS} defined in \eqref{eq:set-gens} is always written $\Ey$, with its argument. We put $\bbK\coloneqq\bbQ(\epsilon_1,\epsilon_2)$ and $\Kal\coloneqq\bbQ(\alpha_0,\alpha_1,\alpha_2,\epsilon_1,\epsilon_2)$. The maximal torus $\Tmax\simeq(\bbC^\times)^5$ acts on $Z$ through an action on $\bbC^6$ with weights $(\alpha_0,\alpha_1,\alpha_2,0,-\epsilon_1,-\epsilon_2)$. It lifts the action on $\bbP^2$ with fixed points $P_0=[1:0:0]$, $P_1=[0:1:0]$, $P_2=[0:0:1]$, whose lift to $\cO_{\bbP^2}(1)$ has weight $-\alpha_i$ at $P_i$. Let $H\coloneqq\mathrm{pr}^*c_1^{\Tmax}(\cO_{\bbP^2}(1))$ be the equivariant line class, so that $H|_{P_i}=-\alpha_i$ and the equivariant cohomology of $Z$ is generated by $H$ with the relation $\prod_{i}(H+\alpha_i)=0$. The elementary symmetric functions of $\alpha=(\alpha_0,\alpha_1,\alpha_2)$ are written $e_k(\alpha)$.

The tangent weights at $P_i$, for $\{i,j,k\}=\{0,1,2\}$, are
\[
\wts_i\coloneqq\{3\alpha_i,\ \alpha_j-\alpha_i,\ \alpha_k-\alpha_i,\ -\epsilon_1,\ -\epsilon_2\},
\qquad \wprod_i\coloneqq\prod_{a\in\wts_i}a .
\]
The equivariant pairing is $\eta(\phi,\phi')\coloneqq\sum_{k}\phi|_{P_k}\phi'|_{P_k}/\wprod_k$. In the basis $\{H^a\}_{a=0,1,2}$ its Gram matrix is
\begin{equation}
\label{eq:set-eta}
\eta_{ab}=\eta(H^a,H^b)=\frac{1}{3\epsilon_1\epsilon_2}
\begin{pmatrix}
\frac{1}{e_3(\alpha)} & 0 & 0\\
0 & 0 & -1\\
0 & -1 & e_1(\alpha)
\end{pmatrix}_{ab},
\qquad \eta^{ab}\coloneqq(\eta^{-1})_{ab}.
\end{equation}
\draftnote{BS label: eq:etaflat.}

The \emph{fibrewise torus} $\TF\subset\Tmax$ is the two-dimensional subtorus given by $\alpha=\alpha^*\coloneqq(E/3)(1,1,1)$. It acts trivially on $\bbP^2$. On $\TF$ we also write $\alpha\coloneqq E/3$. The weights of the normal bundle $\Ntw\coloneqq\cO(-3)_{3\alpha}\oplus\cO_{-\epsilon_1}\oplus\cO_{-\epsilon_2}$ on $\TF$ are recorded by the triple
\[
\bw\coloneqq(E,-\epsilon_1,-\epsilon_2).
\]

\emph{Invariants.} For $g\geq 2$ and $d\geq 0$, the invariants $\GW_{g,d}^{\Tmax}$ are the degrees of the virtual class of $\Mbar_g(Z,d)$ defined by virtual localisation, as in \cite{BS}; the fixed locus is proper by \cref{lem:gen-fixed}. On $\TF$ the same definition reads
\[
\GW_{g,d}^{\TF}=\int_{[\Mbar_g(\bbP^2,d)]^{\vir}}e_{\TF}\big(R\pi_*f^*\Ntw\big)^{-1},
\]
see \cref{prop:gen-poly}. The generating series are
\begin{equation}
\label{eq:set-series}
\sF_g^{\Tmax}\coloneqq\sum_{d\geq 0}\cQ^d\,\GW_{g,d}^{\Tmax},\qquad
\sF_g^{\TF}\coloneqq\sum_{d\geq 0}\cQ^d\,\GW_{g,d}^{\TF}.
\end{equation}
\draftnote{BS: definition of the refined GW series by localisation (labels defn: refined GW series, notation: defn via localisation).}

\begin{rmk}[The term $d=0$]
\label{rmk:set-d0}
The series \eqref{eq:set-series} include the term $d=0$. Throughout the paper $\sF_g$ is read with $d\geq0$, which is the reading under which the identity with the stable graph sum holds on the nose. With a sum over $d>0$ only, as in the definition of the generating series in \cite{BS}, the two series differ by the constant $\GW_{g,0}^{\TF}\in\bbK$, which is computed in \cref{prop:spe-const}.
\end{rmk}

\emph{Correlators and quantum product.} For a point $\ft=\sum_a\ft_aH^a$ of the state space, $\langle\cdots\rangle_{g,n,d}^Z$ denotes the $\Tmax$-equivariant descendent correlators of $Z$ defined by localisation, and $\langle\!\langle\cdots\rangle\!\rangle_{g,n}^Z$ the generating correlators at $\ft$, with the convention fixed before \cref{prop:gen-qde}. The quantum product $\qprod$ is defined by
\begin{equation}
\label{eq:set-qprod}
\eta\big(H^a\qprod H^b,H^c\big)\coloneqq\big\langle\!\big\langle H^a,H^b,H^c\big\rangle\!\big\rangle_{0,3}^Z,
\end{equation}
with $\cQ\coloneqq\re^{\ft_1}$. We write $\eid_i$ for its idempotents, $\Delta_i\coloneqq\eta(\eid_i,\eid_i)^{-1}$ and $\fid_i\coloneqq\Delta_i^{1/2}\eid_i$. The fundamental solution of the Dubrovin connection is
\begin{equation}
\label{eq:set-fundsol}
\Fsol(\ft,z)(H^a)\coloneqq H^a+\sum_{b,c=0}^{2}H^c\,\eta^{cb}\Big\langle\!\!\Big\langle H^a,\frac{H^b}{z-\psi}\Big\rangle\!\!\Big\rangle_{0,2}^Z,
\end{equation}
and on the small locus $\ft_0=\ft_2=0$, by the divisor equation, the small $J$-function is
\begin{equation}
\label{eq:set-J}
\Jsm(\cQ,z)=z\,\cQ^{H/z}\Big(\mathbf{1}+\sum_{a,b}\sum_{d\geq1}\cQ^dH^a\eta^{ab}\Big\langle\frac{H^b}{z(z-\psi_1)}\Big\rangle_{0,1,d}^Z\Big).
\end{equation}
It determines $\Fsol(\mathbf{1})=\Jsm/z$ and $\Fsol(H)=\cQ\partial_{\cQ}\Jsm$ on the small locus. \draftnote{BS labels: eq:qprod, eq:fundsol, eq:Jred, eq:SJ1.}

%-------------------------------------------------------------------------------
\subsection{Mirror data}
\label{sec:set-mirror}

Put $\theta\coloneqq y\,\dd/\dd y$ (the symbol $\theta$ always means $\theta_y$), $\xi\coloneqq1+27y$ and $X\coloneqq1/\xi$. Note that $\theta=(\xi-1)\,\dd/\dd\xi$. The $I$-function of $Z$ is
\begin{equation}
\label{eq:set-I}
I(y,z)\coloneqq z\,y^{H/z}\sum_{d\geq0}y^d\,\frac{\prod_{m=0}^{3d-1}(-3H-mz)}{\prod_{j=0}^{2}\prod_{m=1}^{d}(H+\alpha_j+mz)},
\end{equation}
convergent for $|y|<1/27$. It satisfies the equivariant Picard--Fuchs equation $\LPF I=0$, where
\begin{equation}
\label{eq:set-pf}
\LPF\coloneqq\prod_{j=0}^{2}(z\theta+\alpha_j)+z^3y\prod_{j=0}^{2}(3\theta+j).
\end{equation}
Write $I^{[a]}_n(y)\coloneqq[H^az^{-n}]I(y,z)$, and put $a_d\coloneqq(3d)!/(d!)^3$ and $H_n\coloneqq\sum_{k=1}^n1/k$. Expanding \eqref{eq:set-I},
\begin{equation}
\label{eq:set-Ian}
\begin{gathered}
I_0^{[1]}=\log y+\gI,\qquad \gI\coloneqq3\sum_{d\geq1}\frac{(3d-1)!}{(d!)^3}(-y)^d,\qquad
I_1^{[1]}=-3e_1(\alpha)\sum_{d\geq1}\frac{(3d-1)!\,H_d}{(d!)^3}(-y)^d,\\
I_1^{[2]}=\tfrac12\log^2y+\gI\log y+\gtwo,\qquad \gtwo\coloneqq9\sum_{d\geq1}\frac{(3d-1)!\,(H_{3d-1}-H_d)}{(d!)^3}(-y)^d.
\end{gathered}
\end{equation}
\draftnote{BS labels: eq:Ifun, eq:pf, eq:Ian. BS writes the logarithmic part of $I_1^{[2]}$ as $-\frac12\log^2y+\log y\,I_0^{[1]}$ and $\gtwo$ with digamma values; the two forms agree.}

The \emph{mirror map} is
\[
\cQ(y)\coloneqq\re^{I_0^{[1]}(y)}=y\exp\Big(3\sum_{d\geq1}\frac{(3d-1)!}{(d!)^3}(-y)^d\Big).
\]
It is a unit times $y$ with coefficients in $\bbQ$ and independent of $\alpha$, so that $\Kal\fps{\cQ}=\Kal\fps{y}$. Through it, $\sF_g^{\TF}(y)\coloneqq\sum_{d\geq0}\cQ(y)^d\,\GW_{g,d}^{\TF}$ is an element of $\bbK\fps{y}$ for $g\geq2$. The generators of \cite{BS} are the following functions on $|y|<1/27$:
\begin{equation}
\label{eq:set-gens}
\begin{gathered}
X=\frac{1}{1+27y},\qquad
A\coloneqq\theta\log\cQ=\theta I_0^{[1]}=\sum_{d\geq0}a_d(-y)^d={}_2F_1\big(\tfrac13,\tfrac23;1;-27y\big),\\
D\coloneqq\theta\big(A^{-1}\theta I_1^{[2]}\big),\qquad
\Ey\coloneqq\frac{\theta\big(A^{-1}\theta I_1^{[1]}\big)}{e_1(\alpha)},\qquad
F\coloneqq\theta\log A.
\end{gathered}
\end{equation}
\draftnote{BS label: eq:ADEFXdef.} The letters $\sfF$, $\sfA$ denote indeterminates whose values are $F$, $A$. We put $\thQ\coloneqq\cQ\,\dd/\dd\cQ=A^{-1}\theta$.

Further,
\begin{equation}
\label{eq:set-LW}
\Lhat(Y)\coloneqq\prod_{j=0}^{2}(Y+\alpha_j)+27yY^3,\qquad W(Y)\coloneqq e_1(\alpha)Y^2+2e_2(\alpha)Y+3e_3(\alpha).
\end{equation}
At generic weights, $Y_i$ denotes the root of $\Lhat$ with $Y_i(0)=-\alpha_i$, $W_i\coloneqq W(Y_i)$, and $\mathsf{u}_i$ is a primitive with $\theta\mathsf{u}_i=Y_i$; see \cref{lem:gen-Y}. \draftnote{BS labels: eq:LLMM, eq:delta1.}

The \emph{superpotential} and the primitive form are
\begin{equation}
\label{eq:set-spot}
\cW(y,\alpha,q)\coloneqq\sum_{i=0}^{2}(q_i-\alpha_i\log q_i)+\Big(\frac{\prod_{i=0}^2q_i}{y}\Big)^{1/3},\qquad
\dd\Omega\coloneqq\frac13\prod_{i=0}^{2}\dd\log q_i.
\end{equation}
\draftnote{BS label: eq:spotP2.}

\emph{Formal Laplace expansion.} Let $\phi$ have a non-degenerate critical point $q^{\mathrm{cr}}$ with Hessian $\Hess_{ab}$ and inverse $\Hess^{ab}$, and let $\phi^{\mathrm{ng}}\coloneqq\phi(q)-\phi(q^{\mathrm{cr}})-\frac12\sum_{a,b}\Hess_{ab}(q_a-q_a^{\mathrm{cr}})(q_b-q_b^{\mathrm{cr}})$. Following \cite{BS},
\begin{equation}
\label{eq:set-asym}
\Asym_{q=q^{\mathrm{cr}}}\big[\re^{-\phi/z}\gamma\,\dd^nq\big]\coloneqq(\det\Hess)^{-1/2}\exp\Big[\frac z2\sum_{a,b}\Hess^{ab}\partial_a\partial_b\Big]\big(\re^{-\phi^{\mathrm{ng}}/z}\gamma\big)\Big|_{q=q^{\mathrm{cr}}}.
\end{equation}
The integrand $\re^{\cW/z}\dd\Omega$ of the mirror corresponds to $\phi=-\cW$. \draftnote{BS: the definition of Asym precedes exmpl:stirling.}

%-------------------------------------------------------------------------------
\subsection{Colours on the fibrewise torus}
\label{sec:set-colours}

On $\TF$ put $t\coloneqq y^{1/3}$, $\omega\coloneqq\re^{2\pi\ri/3}$ and, for $i=0,1,2$, $t_i\coloneqq\omega^it$. Then
\begin{equation}
\label{eq:set-colours}
Y_i\coloneqq-\frac{E/3}{1+3t_i},\qquad
q_i\coloneqq Y_i+\frac E3=\frac{Et_i}{1+3t_i},\qquad
v_i\coloneqq\frac{1}{q_i}=\frac{1+3t_i}{Et_i}=\frac{3}{E+3Y_i},\qquad
V_i\coloneqq Ev_i .
\end{equation}
The $Y_i$ are the three roots of $\Lhat$ at $\alpha=\alpha^*$. On $\TF$ the function $W$ of \eqref{eq:set-LW} takes the values
\begin{equation}
\label{eq:set-WTF}
W(Y_i)=E\,q_i^2=\frac E9(3Y_i+E)^2=\frac{E^3t_i^2}{(1+3t_i)^2}.
\end{equation}
We define
\begin{equation}
\label{eq:set-DeltaTF}
\Delta_i\coloneqq3\epsilon_1\epsilon_2W(Y_i)=\frac{3\epsilon_1\epsilon_2E}{v_i^2},\qquad
\Delta_i^{1/2}\coloneqq\frac{\sqrt{3\epsilon_1\epsilon_2E}}{v_i},
\end{equation}
where the second formula fixes the square root. It is transported together with the rows under $t\mapsto\omega^it$. At generic weights $\Delta_i=3\epsilon_1\epsilon_2W_i$ is a statement proved in \cref{lem:gen-Y}; here \eqref{eq:set-DeltaTF} is a definition.

Since $\xi=(1+3t)(1-3t+9t^2)$, and $1+3\omega^it=1-\omega^i\neq0$ at $t=-1/3$ for $i=1,2$, colour $0$ on the real branch $t\to-1/3$ is the unique colour with $Y_i\to\infty$ as $y\to-1/27$. We call it the \emph{singular colour}, and from now on $t$ is on this branch near $\xi=0$, that is $t=-(1-\xi)^{1/3}/3$ for $\xi$ small.

\emph{Row and edge variables.} The letter $z$ denotes the argument of a row. The letters $\psi,\psi',w$ denote edge variables: a half-edge with variable $\psi$ carries the row at $z=-\psi$.

%-------------------------------------------------------------------------------
\subsection{Leaves}
\label{sec:set-leaves}

For $j\in\bbZ\setminus\{0\}$ put $\lf_j(\psi)\coloneqq1/(E-j\psi)$, and put $\lf_0\coloneqq1/E$. Further, $\Lsp_m\coloneqq\mathrm{span}\{\lf_1,\ldots,\lf_m\}$, $\Lsp_0\coloneqq0$, $\Lsp_m^{+}\coloneqq\bbC+\Lsp_m$; in bounds, $\Lsp_m$ is read as $\Lsp_{\max(m,0)}$. The following identities are elementary:
\begin{align}
-\psi\,\lf_j&=\frac{1-E\lf_j}{j}, && j\neq0, \label{eq:set-lam1}\\
(E-N\psi)\,\lf_j&=\frac Nj+\Big(1-\frac Nj\Big)E\lf_j, && j\neq0;\quad\text{in particular } (E-N\psi)\lf_N=1 \text{ for } N\neq0, \label{eq:set-lam2}\\
\lf_l(\psi)-\lf_l(-w)&=l(\psi+w)\,\lf_l(\psi)\,\lf_l(-w), \label{eq:set-lam3}\\
E\lf_1&=1+\psi\lf_1. \label{eq:set-lam4}
\end{align}
The base point and the dilaton tail are
\[
\qo\coloneqq E-E^2\lf_1=-\frac{E\psi}{E-\psi},\qquad \Tdil\coloneqq\qo+\psi=-\frac{\psi^2}{E-\psi}.
\]

%-------------------------------------------------------------------------------
\subsection{The conifold coordinate and the unit \texorpdfstring{$\rho$}{rho}}
\label{sec:set-coord}

Let
\[
\Ltwo\coloneqq\theta^2+3y(3\theta+1)(3\theta+2)=\xi\theta^2+27y\theta+6y.
\]
Its exponents at $\xi=0$ are $(0,0)$, because near $\xi=0$ it reads $\xi\partial_\xi^2+\partial_\xi+\cdots$. Let $\wcf$ be the solution analytic at $\xi=0$ with $\wcf(0)=1$, and let $\tc$ be defined by
\[
\theta\tc=-\wcf,\qquad\tc(\xi=0)=0.
\]
Put $\ccf\coloneqq\theta\tc=-\wcf$. For any non-vanishing section $\cone(y)$ put
\[
\Fc\coloneqq\theta\log\cone,\qquad\dtwo\coloneqq\frac{E(\cone^{-2}-1)}{2\xi}.
\]

\begin{lem}\genusstatus{AI-proposed, unchecked}
\label{lem:set-coord}
\begin{description}
\item[(a)] $\tc=\xi+O(\xi^2)$ is an analytic coordinate at $\xi=0$, and $\ccf(0)=-1$. At $\cone=\ccf$, the functions $\cone^{-1}$, $\Fc$ and $\dtwo$ are analytic at $\xi=0$; for $\dtwo$ this holds because $\cone(0)^{-2}=1$.
\item[(b)] On the real branch $t=-(1-\xi)^{1/3}/3$ one has $v_0=-\tc/E+O(\tc^2)$. Hence
\[
\rho\coloneqq-\frac{Ev_0}{\tc}
\]
is an analytic unit with $\rho(0)=1$, and $\rho^2\Delta_0=3\epsilon_1\epsilon_2E^3/\tc^2$.
\end{description}
\end{lem}

\begin{proof}
(a) From the definition and $\theta=(\xi-1)\dd/\dd\xi$ one gets $\dd\tc/\dd\xi=-\wcf/(\xi-1)$, which equals $1$ at $\xi=0$. The remaining assertions follow from $\ccf=-\wcf$ and $\wcf(0)=1$.

(b) We have $\xi=(1+3t)(1-3t+9t^2)$ and the second factor equals $3$ at $t=-1/3$. So $1+3t=\xi/3+O(\xi^2)$ and $v_0=(1+3t)/(Et)=-\xi/E+O(\xi^2)$. With (a) this gives the expansion of $v_0$ and the unit $\rho$. The identity follows from $\Delta_0=3\epsilon_1\epsilon_2E/v_0^2$.
\end{proof}

The comparison of $\tc$ with the coordinate $\tcf$ of \cite{BS} is \cref{lem:fra-tcf}.

\begin{rmk}[Normalisation of $\tc$]
\label{rmk:set-norm}
The normalisation of $\tc$ enters every exact constant through $\rho(0)=1$, and through the identity \eqref{eq:row-zinf}. The filtration statements (the memberships of \cref{thm:lea-main}) are unchanged if $\tc$ is rescaled as an expansion variable, $\tc\mapsto a\tc$ with $\cone=\ccf$ kept: this multiplies $\rho$ by the constant $1/a$, so that $\rho(0)=1/a$ and the base point is $\qo/a$. The domain $\Dom$ of \cref{sec:set-hodge} and the homogeneity of \cref{lem:reg-H0}, as written, need $\rho(0)=1$. Such a rescaling is not a change of $\cone$: for $a^2\neq1$ the section $\cone=a\theta\tc$ gives $\cone(0)^{-2}\neq1$, and then $\dtwo$ acquires a pole. The statements are also not stable under a non-linear change of coordinate with the same leading term: with $\xi$ in place of $\tc$ in the definition of $\rho$, the constant term of $[\tc^1](\rho\,\qt)$ is $(11/18)E\neq0$, because \eqref{eq:row-zinf} is the exact identity $\Sinf=-\tc$.
\end{rmk}

%-------------------------------------------------------------------------------
\subsection{Special functions}
\label{sec:set-special}

Let
\[
\Theta(x)\coloneqq\exp\sum_{m\geq1}\frac{B_{2m}\,x^{2m-1}}{2m(2m-1)},\qquad\text{so that }\Theta(-x)=\frac{1}{\Theta(x)}.
\]
The two-variable version is that of \cite{BS}, with $B_n(x)$ the Bernoulli polynomials:
\[
\Theta(z,x)\coloneqq\exp\Big[-\sum_{k\geq1}\frac{B_{k+1}(x)(-z)^k}{k(k+1)}\Big].
\]
Then $\Theta(z,0)=\Theta(z)$, $\Theta(z,x+1)=(1+xz)\Theta(z,x)$ and $\Theta(z,-m)=\Theta(z)/\prod_{j=1}^{m}(1-jz)$ for integers $m\geq0$. For a multiset $\mathsf{v}$ of non-zero weights put
\[
h_{\mathsf{v}}(z)\coloneqq\prod_{a\in\mathsf{v}}\Theta(-z/a).
\]
In particular
\begin{equation}
\label{eq:set-hw}
\hb(z)=\Theta(-z/E)\,\Theta(z/\epsilon_1)\,\Theta(z/\epsilon_2),\qquad\hb(z)\hb(-z)=1,
\end{equation}
which is the series $\Upsilon(z/\epsilon_1,z/\epsilon_2)$ of \cite{BS}. We also put
\[
\Nz(z)\coloneqq\sqrt{\epsilon_1\epsilon_2}\,\Theta(-z/\epsilon_1)\Theta(-z/\epsilon_2).
\]
\draftnote{BS: Theta and Upsilon are defined in exmpl:stirling.}

The numbers $r_l$ are defined by the recursion
\begin{equation}
\label{eq:set-rec}
r_{-1}=0,\quad r_0=1,\qquad 27(m+1)\,r_{m+1}=3(3m^2+3m+1)\,r_m-m^3\,r_{m-1}\quad(m\geq0),
\end{equation}
so that $r_0,r_1,r_2,r_3=1,\frac19,\frac2{81},\frac{14}{2187}$. The closed-form series are
\begin{equation}
\label{eq:set-S}
\Srow(v,z)\coloneqq\sum_{l\geq0}\frac{r_l\,(vz)^l}{\prod_{j=1}^{l+1}(1+jz/E)},\qquad
\Srow^{[m]}(v,z)\coloneqq\sum_{l\geq0}\frac{r_l\,(vz)^l}{\prod_{j=1}^{l+1-m}(1+jz/E)},
\end{equation}
where a product with upper limit $<1$ is read as $1$. The case $l=0$, $m=2$ occurs in \cref{lem:row-closed12}. Each $z$-coefficient of $\Srow$ is a polynomial in $v$, and $[z^k]\Srow$ has degree $\leq k$ in $v$: this is immediate from the definition, since only $l\leq k$ contribute.

%-------------------------------------------------------------------------------
\subsection{The rows}
\label{sec:set-rows}

\begin{defn}[Frame-B rows]
\label{def:set-rows}
At colour $i$ and for a non-vanishing section $\cone(y)$, put $\Thp_i\coloneqq Y_i+z\theta$ and
\begin{align*}
\PB_0^{(i)}&\coloneqq\Delta_i^{-1/2}\,\Srow(v_i,z),\\
\PB_1^{(i)}&\coloneqq\cone^{-1}\,\Thp_i\,\PB_0^{(i)},\\
\PB_2^{(i)}&\coloneqq\xi\,\cone^{2}\,\big[\Thp_i\,\PB_1^{(i)}-\dtwo\,\PB_1^{(i)}\big].
\end{align*}
The \emph{large-radius-frame rows} are $\RLR_a^{(i)}\coloneqq\hb(z)\,\PB_a^{(i)}$. The \emph{normalised rows} are $\Ph_a^{(i)}\coloneqq\Delta_i^{1/2}\PB_a^{(i)}$. Thus $\Ph_0^{(i)}=\Srow(v_i,z)$ is component $0$ of the colour-$i$ normalised row $\Ph^{(i)}=(\Ph_0^{(i)},\Ph_1^{(i)},\Ph_2^{(i)})$, not the whole row. The main theorems take $\cone=\ccf$, the same section at every colour; this is legitimate because $\cone$ depends on $y$ only.
\end{defn}

\begin{rmk}
\begin{enumerate}
\item The definition uses only $\Fc=\theta\log\cone$, through $\theta\cone^{-1}=-\cone^{-1}\Fc$, so $\Ltwo$ plays no role in defining the rows.
\item With $\cone=A$ these are the rows of the row recursion of \cite{BS}, through $\xi A^2=1/D$ and $\dtwo=e_1(\alpha^*)\Ey$; both are \cref{prop:row-modrel}, see \cref{rmk:row-dict}. The case $\cone=\ccf$ is the conifold substitution of \cref{def:fra-cf} applied to these rows, see \cref{sec:fra}. \draftnote{BS label: eq:RHb.}
\end{enumerate}
\end{rmk}

\emph{Pairing on the rows.} For vectors with components indexed by $a,b\in\{0,1,2\}$, put $\langle\mathsf{a},\mathsf{b}\rangle\coloneqq\sum_{a,b}\eta^{ab}\mathsf{a}_a\mathsf{b}_b$, with
\begin{equation}
\label{eq:set-etaTF}
\eta^{00}=\frac{\epsilon_1\epsilon_2E^3}{9},\qquad\eta^{11}=-3\epsilon_1\epsilon_2E,\qquad\eta^{12}=\eta^{21}=-3\epsilon_1\epsilon_2,
\end{equation}
and all others $0$. Note that $\eta^{11}=E\eta^{12}$ and $\eta^{22}=0$. These are the values of the inverse of \eqref{eq:set-eta} at $\alpha=\alpha^*$. In \cref{sec:row,sec:lea,sec:reg}, as far as the decoration sum is concerned, \eqref{eq:set-etaTF} is a definition.

%-------------------------------------------------------------------------------
\subsection{The factors}
\label{sec:set-factors}

The \emph{vertex weight} at colour $i$ is $\Delta_i$.

The \emph{dilaton factor} is
\begin{equation}
\label{eq:set-dil}
\Dilh_i(d)\coloneqq(-1)^d[z^{d-1}]\Ph_0^{(i)}\quad(d\geq1),\qquad\text{equivalently}\qquad
\qt^{(i)}(\psi)\coloneqq\sum_{d\geq1}\Dilh_i(d)\psi^d=-\psi\,\Srow(v_i,-\psi).
\end{equation}
Since $[z^0]\Srow=1$, we have $\Dilh_i(1)=-1$ for every $t$: the term $d=1$ is the dilaton shift, and the decorations below have $d\geq2$. The factor $\Dilh_i$ uses only component $0$ of the row, so it does not depend on $\cone$. We write $\qt\coloneqq\qt^{(0)}$.

The \emph{edge factor}, with edge variables $\psi,w$ (rows at $z=-\psi$ and $z=-w$), is defined by
\begin{equation}
\label{eq:set-edge}
\begin{gathered}
\cP_{ij}(\psi,w)\coloneqq-\big\langle\Ph^{(i)}(-\psi),\Ph^{(j)}(-w)\big\rangle,\
\Edgeh_{ij}(\psi,w)\coloneqq\frac{\cP_{ij}(\psi,w)-\cP_{ij}(-w,w)}{\psi+w},\qquad
\Edgeh_{ij}(m,n)\coloneqq[\psi^mw^n]\Edgeh_{ij}.
\end{gathered}
\end{equation}
Equation \eqref{eq:set-edge} is a definition. Its form as the Givental edge in $\Delta$-normalisation, and the symmetry $\Edgeh_{ij}(\psi,w)=\Edgeh_{ji}(w,\psi)$, are stated in \eqref{eq:row-edgedelta}, after the symplecticity \cref{prop:row-sympl}.

\begin{rmk}[Signs of square roots]
\label{rmk:set-signs}
In the normalised form \eqref{eq:set-graphsum} below no square root occurs. Indeed, $\Delta_i=3\epsilon_1\epsilon_2E/v_i^2$, and $\Ph_a=\Delta_i^{1/2}\PB_a$ does not depend on the sign of $\Delta_i^{1/2}$ because $\PB_0$ carries $\Delta_i^{-1/2}$. So $\Ph_a$ is built from $\Srow$, $\theta$, $Y_i$, $\cone^{-1}$, $\dtwo$ and $\theta\log\Delta_i$, all single valued. Hence each $\GCont(\Gamma,\col)$ below is well defined on any cover. The sign of $\cone$ is a different matter: see \cref{rmk:lea-gauge}.
\end{rmk}

%-------------------------------------------------------------------------------
\subsection{The Hodge class and the decoration sum}
\label{sec:set-graphsum}

For a stable curve of genus $h$, with $\lambda_i=c_i(\Hdg)$, let $\Lamv(u)\coloneqq\sum_i(-1)^i\lambda_iu^{h-i}$. For a multiset $\mathsf{v}$ of non-zero weights, the \emph{normalised Hodge class} is $\Hg_{\mathsf{v}}\coloneqq\prod_{a\in\mathsf{v}}\Lamv(a)/a^h$. We use
\begin{equation}
\label{eq:set-Hg}
\Hg_{\bw}=\prod_{a=1}^{3}\frac{\Lamv(\bw_a)}{\bw_a^{h}},\qquad\bw=(E,-\epsilon_1,-\epsilon_2).
\end{equation}
It is the same at every colour and independent of $y$. The scale law of \cref{cor:van-scale} depends on this normalisation: with the unnormalised product the exponent there would change. The class is a pull-back under forgetful maps, so the string and dilaton equations hold for $\Hg_{\bw}$ when $h\geq2$.

\begin{defn}[Decoration sum]
\label{def:set-graphsum}
Fix $g\geq2$. Let $\Gamma$ run over connected stable graphs of genus $g$ without legs, with vertex set $\Ver(\Gamma)$, edge set $\Edg(\Gamma)$, vertex genera $g_v$, and $n_v$ half-edges at $v$. A \emph{colouring} is a map $\col:\Ver(\Gamma)\to\{0,1,2\}$, $v\mapsto\col_v$. A \emph{decoration} assigns an exponent $a_h\geq0$ to each half-edge $h$ (we write $a_e,b_e$ for the exponents at the two half-edges of an edge $e$) and to each vertex $v$ an integer $m_v\geq0$ with an ordered tuple $(d_{v,1},\ldots,d_{v,m_v})$ of \emph{dilaton degrees} $d_{v,k}\geq2$, with weight $1/m_v!$. The \emph{decoration sum} of $(\Gamma,\col)$, with all factors taken at $\cone=\ccf$, is
\begin{multline}
\label{eq:set-graphsum}
\GCont(\Gamma,\col)\coloneqq\sum_{\mathrm{dec}}\ \prod_{v\in\Ver(\Gamma)}\Big[\Delta_{\col_v}^{g_v-1}\frac{1}{m_v!}\int_{\Mbar_{g_v,n_v+m_v}}\Hg_{\bw}\prod_{h\text{ at }v}\psi_h^{a_h}\prod_{k=1}^{m_v}\psi_{n_v+k}^{d_{v,k}}\ \prod_{k=1}^{m_v}\Dilh_{\col_v}(d_{v,k})\Big]\\
\times\prod_{e\in\Edg(\Gamma)}\Edgeh_{\col_{v(e)}\col_{v'(e)}}(a_e,b_e),
\end{multline}
where $v(e),v'(e)$ are the end points of $e$. The \emph{frame-B Hodge-vertex graph sum} is
\begin{equation}
\label{eq:set-Gsum}
\Gsum_g\coloneqq\sum_{\Gamma}\frac{1}{|\Aut\Gamma|}\sum_{\col:\Ver(\Gamma)\to\{0,1,2\}}\GCont(\Gamma,\col).
\end{equation}
\end{defn}

The decoration sum is finite. Since $\Hg_{\bw}$ has mixed degree, the integral vanishes unless $\sum a_h+\sum d_{v,k}\leq3g_v-3+n_v+m_v$. As $d_{v,k}\geq2$, this forces $m_v\leq3g_v-3+n_v$. So $\GCont(\Gamma,\col)$ is a finite sum of products of the factors above, and this is what gives ``analytic at the conifold'' its meaning in \cref{thm:reg-main}. By the edge symmetry \eqref{eq:row-edgedelta}, $\GCont(\Gamma,\col)$ does not depend on the orientation of the edges.

Equation \eqref{eq:set-graphsum} is the definition of the sum. The regularity statements of the paper are statements about it. Its identification with a Gromov--Witten potential is a separate statement, \cref{prop:fra-frame}.

%-------------------------------------------------------------------------------
\subsection{The Hodge side}
\label{sec:set-hodge}

Fix $h\geq2$. Let $\RRing\coloneqq\RRing_0\fps{\tc}$, with $\RRing_0$ a $\bbQ$-algebra in which $\epsilon_1,\epsilon_2,E$ are invertible (for instance $\bbC$ or $\bbQ(\omega,\epsilon_1,\epsilon_2)$).

\begin{defn}
\label{def:set-fpot}
\begin{description}
\item[Base-point derivatives.] For $\mathsf{b}_1,\ldots,\mathsf{b}_s\in\RRing\fps{\psi}$, put
\[
\dbase^s(\mathsf{b}_1,\ldots,\mathsf{b}_s)\coloneqq\sum_{m\geq 0}\frac{1}{m!}\int_{\Mbar_{h,s+m}}\Hg_{\bw}\prod_{i=1}^s\mathsf{b}_i(\psi_i)\prod_{k=1}^{m}\Tdil(\psi_{s+k}).
\]
This is a finite sum, since $\Tdil=O(\psi^2)$ forces $m\leq3h-3+s$, and each power series is cut off by the dimension.
\item[The potential on the neighbourhood $\Dom$ of $\qo$.] Let $\Dom\coloneqq\qo+\tc\RRing\fps{\psi}$. For $\mathbf{q}=\qo+\delta\in\Dom$ and $f_1,\ldots,f_n\in\RRing\fps{\psi}$, define
\[
\dd^n\Fpot_h(\mathbf{q})[f_1,\ldots,f_n]\coloneqq\sum_{r\geq0}\frac{1}{r!}\,\dbase^{n+r}(f_1,\ldots,f_n,\underbrace{\delta,\ldots,\delta}_{r}),\qquad\Fpot_h(\mathbf{q})\coloneqq\dd^0\Fpot_h(\mathbf{q}),
\]
which converges $\tc$-adically because $\delta=O(\tc)$. Note that $\delta$ may have a term linear in $\psi$: the domain $\Dom$ contains $\rho\,\qt$, whose $\psi$-coefficient is $-\rho$. So $\mathbf{q}$ need not be of the form $-\psi+O(\psi^2)$.
\item[The dilaton-shifted form.] If $\mathbf{q}=-\psi+O(\psi^2)$, put
\[
\Fpot_h(\mathbf{q})\coloneqq\sum_{m\geq 0}\frac{1}{m!}\int_{\Mbar_{h,m}}\Hg_{\bw}\prod_{i=1}^m(\mathbf{q}+\psi)(\psi_i),
\]
a finite sum, and similarly for $\dd^n\Fpot_h(\mathbf{q})[f_1,\ldots,f_n]$. On the overlap with $\Dom$ the two definitions agree: expand $\prod(\Tdil+\delta)$ binomially, and the factor $\frac{1}{m!}\binom mr$ reindexes to $\frac{1}{r!}\frac{1}{(m-r)!}$. This form is what occurs at the regular colours, where $\qt^{(j)}$ is not in $\Dom$ and no rescaling or Taylor expansion is applied.
\item[The numbers $\Xiz_h$.] For a multiset $\msU=\{j_1,\ldots,j_r\}$ of integers,
\[
\Xiz_h(\msU)\coloneqq\dd^r\Fpot_h(\qo)[\lf_{j_1},\ldots,\lf_{j_r}]=\dbase^r(\lf_{j_1},\ldots,\lf_{j_r})=\sum_{m\geq 0}\frac{1}{m!}\int_{\Mbar_{h,r+m}}\Hg_{\bw}\prod_{i=1}^{r}\lf_{j_i}(\psi_i)\prod_{k=1}^{m}\Tdil(\psi_{r+k}).
\]
Further, $J_{\msU}\coloneqq\sum_ij_i$, $\dg(\msU)\coloneqq\sum_i(j_i-1)$ and $\ZZ[h,n]\coloneqq\Xiz_h(1^n)$, so that $\ZZ[h,0]=\Fpot_h(\qo)$.
\end{description}
\end{defn}

\begin{defn}[The statement $\Mirr(h)$]
\label{def:set-mirr}
$\Mirr(h)$ is the statement: $\Xiz_h(\msM)=0$ for every \emph{non-empty} multiset $\msM$ of integers $\geq2$ with $\dg(\msM)\leq2h-3$.
\end{defn}

The empty multiset is excluded because $\Xiz_h(\emptyset)=\ZZ[h,0]$ is the gap. By \cref{thm:val-main} it is a non-zero element of $\bbK$; this does not specialise, see \cref{rmk:reg-rank}, and no proof step uses it. Parts $0$ and $1$ are not in $\Mirr$: part $1$ is dealt with by \cref{lem:reg-string}, and part $0$ never occurs, see \cref{rmk:reg-nol0}.

\emph{The vertex factor.} By the definition of $\Fpot_h$, and because the decorations of \eqref{eq:set-graphsum} have $d\geq2$ with weight $1/m!$, the sum over dilaton decorations at a vertex $v$ of genus $h\geq2$ and colour $i$, with half-edge series $f_1,\ldots,f_n$ (where $f_r(\psi)=\sum_a(\text{edge coefficient})\,\psi^a$), equals
\begin{equation}
\label{eq:set-vertex}
\Delta_i^{h-1}\,\dd^n\Fpot_h\big(\qt^{(i)}\big)[f_1,\ldots,f_n].
\end{equation}
At colour $0$, $\qt$ lies in $\Dom$ by \cref{lem:lea-D0}. Edges do not factorise in general: a loop, or an edge between two rescaled vertices, is a tensor. The expression \eqref{eq:set-vertex} is multilinear in the half-edge slots, and it is applied term by term, leaf by leaf in the expansions of \cref{thm:lea-E00} and \cref{thm:lea-M}, and coefficient by coefficient in the edge variables.

\begin{rmk}[Frame and weights]
\label{rmk:set-frame}
\emph{Frame B} means the $\Rm$-matrix $\RB=\Rm_Z\,\hb^{-1}$. The Hodge weights are $\bw=(E,-\epsilon_1,-\epsilon_2)$ and not $-\bw$. The transfer of the vanishing statement, \cref{lem:reg-A}, holds for this choice only: there the leaf scale $E$ produced by the value $\Srow(0,z)=E/(E+z)$ equals the weight $\bw_1$. For $\bw\mapsto-\bw$ the leaf scale is not a weight.
\end{rmk}

%-------------------------------------------------------------------------------
\subsection{The Givental action}
\label{sec:set-givental}

We recall the action of \cite{PPZ15} in the form used in \cref{sec:gen,sec:fra}; see \cite[Definitions~2.2, 2.5, 2.13 and Proposition~2.14]{PPZ15}. Let $(V,\eta,\mathbf{1})$ be a state space with unit and let $\Rm(z)\in1+z\End(V)\fps{z}$ be symplectic, that is, $\Rm(z)\Rm^*(-z)=1$ \cite[(5)]{PPZ15}. On a CohFT $\Omega$, the action $\Rm.\Omega$ is the composite of a translation and a graph sum:
\begin{description}
\item[Translation.] $T_{\Rm}(z)\coloneqq z\big(1-\Rm^{-1}(z)\mathbf{1}\big)=O(z^2)$, and for a series $b$, $(T_{b}\Omega)_{g,n}(x_1,\ldots,x_n)\coloneqq\sum_{m\geq0}\frac{1}{m!}\,p_{m*}\,\Omega_{g,n+m}\big(x_1,\ldots,x_n,b(\psi_{n+1}),\ldots,b(\psi_{n+m})\big)$, with $p_m$ the forgetful map.
\item[Legs.] A leg $l$ with insertion $x_l$ carries $\Rm^{-1}(\psi_l)x_l$.
\item[Edges.] An edge with half-edge variables $\psi,\psi'$ carries the bivector
\[
\frac{\eta^{-1}-\big(\Rm^{-1}(\psi)\otimes\Rm^{-1}(\psi')\big)\eta^{-1}}{\psi+\psi'}.
\]
\end{description}
\draftnote{These formulas are written from the way they are used in the proofs of the Mumford relation, of the flatness of $\Rm$ and of the unfolding lemma; to be compared with the print of \cite[Definitions~2.2, 2.5, 2.13]{PPZ15} before submission, including the sign in (5).}

For a semisimple theory with idempotents $\eid_i$ and $\fid_i=\Delta_i^{1/2}\eid_i$, the rows of $\Rm$ are defined by $\Rm_{aj}(z)\coloneqq\eta\big(H^a,\Rm(z)\fid_j\big)$. By symplecticity this is equivalent to $\fid_j^*\Rm^{-1}(z)H^a=\Rm_{aj}(-z)$. So a leg with insertion $H^a$ at colour $j$ carries $\Rm_{aj}(-\psi)$, and
\begin{align*}
T_{\Rm}(z)&=z\sum_i\big(\Delta_i^{-1/2}-\Rm_{0i}(-z)\big)\fid_i=\sum_i\eid_i\,z\big(1-\Delta_i^{1/2}\Rm_{0i}(-z)\big).
\end{align*}
The coefficient of the bivector on $\fid_i\otimes\fid_j$ is $\big[\delta_{ij}-\sum_{a,b}\Rm_{ai}(-\psi)\eta^{ab}\Rm_{bj}(-w)\big]/(\psi+w)$; on $\eid_i\otimes\eid_j$ it is multiplied by $\Delta_i^{1/2}\Delta_j^{1/2}$, and this is how \eqref{eq:set-edge} arises, see \cref{lem:fra-unfold}. The coefficient of $\eid_i$ in $T_{\RB}$ is written $T_{\RB}^{(i)}$. For a scalar series $h$, $\Rm_{h}$ denotes the $\Rm$-action by the scalar $h$.
